\chapter{Второй выход: как мозг управляет химией тела}
\chaptermark{Химический выход}
\label{sec:chemoutput}

\section{Быстрый и медленный выходы}

В главе~\ref{sec:muscles} мы видели быстрый выход машины --- мышцы. Но у неё есть и второй, медленный. Мозг держит в нужных пределах температуру тела, частоту сердца, количество воды и сахара, готовит тело к бегу или ко сну. Для этого он не двигает суставы, а выпускает химические сигналы. Путей два. Первый --- нервы, идущие прямо к внутренним органам: к сердцу, сосудам, кишечнику, железам. Второй --- кровь: некоторые нейроны и железы выбрасывают вещество не в щель к соседу, а прямо в кровоток.

Кровь --- самая широковещательная шина из всех, что встречались в книге. Широковещательные каналы глав~\ref{sec:serocomputer}--\ref{sec:acet} рассылали сигнал на большие области мозга; кровь рассылает его каждой клетке тела. Кровь обходит тело примерно за минуту, так что сообщение доходит до всех за десятки секунд и держится минуты и часы, пока вещество не разрушится. Адреса у такой рассылки нет вовсе. Как и в утверждении~\ref{prop:receptor}, смысл сообщения выбирает получатель: откликаются только клетки, у которых есть рецептор к этому веществу.

\section{Газ и тормоз: два провода к органам}

Лучший пример --- сердце. В сердце есть собственный ритмоводитель, как у \nt{СЕРО}-нейрона главы~\ref{sec:serocomputer}: он сам задаёт ритм, и если отключить все нервы, сердце бьётся примерно сто раз в минуту. Мозг ритм не порождает, а только подстраивает, причём двумя проводами противоположного знака (рис.~\ref{fig:gasbrake}). По одному идёт \nt{НОРА}: это \emph{газ}, он ускоряет ритм. По другому --- \nt{АЦЕТ}: это \emph{тормоз}, он замедляет. Грубо,
\begin{empheq}[box=\keybox]{equation}
f \;\approx\; f_0 \;+\; a_S\,S \;-\; a_P\,P ,
\label{eq:gasbrake}
\end{empheq}
где $f_0\approx100$ ударов в минуту --- собственный ритм ритмоводителя, $S$ и $P$ --- активность газа и тормоза. В покое тормоз нажат, и сердце бьётся шестьдесят-семьдесят раз в минуту; при нагрузке тормоз отпускают и нажимают газ.

\begin{figure}[t]
\centering
\begin{tikzpicture}[>=Stealth,
  blk/.style={draw,very thick,rounded corners=2pt,align=center,font=\small},
  pace/.style={circle,draw,very thick,double,double distance=1.2pt,minimum size=1.8cm,inner sep=0pt,font=\scriptsize,fill=red!8,align=center},
  inh/.style={-{Bar[width=7pt]},thick,red!70!black}]
\node[blk,fill=blue!5,text width=1.6cm] (B) at (0,0) {команда\\мозга};
\node[pace] (H) at (6.2,0) {ритмо-\\водитель};
\draw[->,very thick,orange!85!black] (B.north east) to[bend left=20] node[above,font=\scriptsize,black,align=center] {\nt{НОРА}: газ, секунды} (H.north west);
\draw[inh,very thick,blue!60!black] (B.south east) to[bend right=20] node[below,font=\scriptsize,black,align=center] {\nt{АЦЕТ}: тормоз, доли секунды} (H.south west);
\draw[->,very thick] (H) -- (8.4,0) node[right,font=\small] {частота $f$};
\node[blk,fill=orange!10,text width=1.6cm] (A) at (0,-2.6) {\nt{АДРЕ}\\в кровь};
\draw[->,thick,orange!85!black] (B) -- (A);
\foreach \y/\lab in {-2.0/{сердце},-2.6/{сосуды},-3.2/{печень, мышцы}}{
  \draw[->,thick,densely dashed,orange!85!black] (A.east) -- (4.3,\y) node[right,font=\scriptsize,black] {\lab};}
\end{tikzpicture}
\caption{Два провода противоположного знака к ритмоводителю сердца: газ \nt{НОРА} медленный, тормоз \nt{АЦЕТ} быстрый. Внизу --- тот же газ по широковещательной шине: \nt{АДРЕ}-клетки выбрасывают адреналин в кровь, и сигнал получают все органы с подходящим рецептором (штриховые стрелки).}
\label{fig:gasbrake}
\end{figure}

Это двухпроводный код главы~\ref{sec:glutcomputer} и пара мышц главы~\ref{sec:muscles}, только теперь провода управляют не суставом, а темпом ритмоводителя. И у проводов разная скорость. Оба работают как фильтры нижних частот, но тормоз пропускает перемены в несколько раз быстрее газа~\cite{Berger1989}: \nt{АЦЕТ} открывает ионные каналы прямо, а \nt{НОРА} действует через цепочку реакций внутри клетки. Инженер узнает здесь приём с двумя приводами разной скорости: быстрый делает точную мгновенную подстройку, медленный --- крупные долгие перемены.

Здесь же находит роль \nt{АДРЕ}, про который в таблице~\ref{tab:types} было написано, что его роль в мозге неясна. Вне мозга она ясна. Часть клеток газового пути не посылает провод к органу, а выбрасывает адреналин прямо в кровь. Это тот же газ, но по широковещательной шине: за десятки секунд он доходит до сердца, сосудов, печени и мышц сразу и держится минуты. Адресный газ \nt{НОРА} подстраивает один орган; широковещательный газ \nt{АДРЕ} переводит в режим бега всё тело.

\section{Каскад с обратной связью}

Многие гормоны выпускаются не одной железой, а каскадом. Небольшая группа нейронов в основании мозга выбрасывает в кровь первый сигнал; он заставляет железу-посредника выпустить второй; тот --- конечную железу выпустить гормон, который и действует на тело. А конечный гормон возвращается к первым ступеням и тормозит их. Получается усилитель из трёх каскадов, охваченный отрицательной обратной связью, --- регулятор уровня, как у \nt{СЕРО}-системы в формуле~\eqref{eq:feedback}.

Опишем каждую ступень RC-цепью с постоянной времени $\tau$ и усилением $g$, а обратную связь --- коэффициентом $k$:
\begin{equation}
\tau\,\frac{\dd x_1}{\dd t} = -x_1 + g\bigl[u - k\,x_3\bigr]_+ ,\qquad
\tau\,\frac{\dd x_2}{\dd t} = -x_2 + g\,x_1 ,\qquad
\tau\,\frac{\dd x_3}{\dd t} = -x_3 + g\,x_2 ,
\label{eq:cascade}
\end{equation}
где $u$ --- команда мозга, $x_3$ --- уровень конечного гормона. Усиление петли $K=g^3k$. В равновесии $x_3=g^3u/(1+K)$, и, как у любого регулятора с обратной связью, большая петля делает уровень нечувствительным к разбросу усилений ступеней. Но каждая ступень сдвигает фазу: на частоте $\omega=\sqrt3/\tau$ три одинаковые ступени вместе дают сдвиг в полоборота и ослабление в восемь раз. Отсюда классический результат теории регулирования:
\begin{empheq}[box=\keybox]{equation}
K \;<\; 8 ,
\label{eq:cascadestable}
\end{empheq}
иначе регулятор начинает колебаться, с периодом около $2\pi\tau/\sqrt3\approx3.6\,\tau$.

\begin{keyblock}
\begin{proposition}[Точность против устойчивости]
\label{prop:cascade}
У трёхступенчатого каскада с пропорциональной обратной связью ошибка уровня равна $1/(1+K)$, а устойчив он только при $K<8$. Поэтому точнее, чем примерно на десять процентов, такой регулятор уровень не держит; попытка поднять усиление превращает его в генератор колебаний.
\end{proposition}
\end{keyblock}

Мы проверили это на модели~\eqref{eq:cascade} (рис.~\ref{fig:cascade}). При $K=4$ уровень после включения немного колеблется и устанавливается. При $K=7$ он тоже устанавливается, но медленно. При $K=12$ колебания не затухают и не растут без предела: ступень не может выдать отрицательный сигнал, и каскад выходит на ровные пульсации между $0.83$ и $1.24$ от среднего уровня с периодом $3.7$--$3.9\,\tau$.

\begin{figure}[t]
\centering
\begin{tikzpicture}[>=Stealth,x=0.3cm,y=1.2cm]
\draw[->,gray!70!black] (0,0) -- (41.5,0) node[right,font=\small,black] {$t/\tau$};
\draw[->,gray!70!black] (0,0) -- (0,2.8) node[left,font=\small,black] {$x_3/x_3^*$};
\foreach \t in {0,10,20,30,40}{\draw[gray!60!black] (\t,-0.05) -- (\t,0.05); \node[below,font=\scriptsize] at (\t,0) {\t};}
\foreach \f in {1,2}{\draw[gray!60!black] (-0.3,\f) -- (0.3,\f); \node[left,font=\scriptsize] at (0,\f) {\f};}
\draw[densely dashed,gray!60!black] (0,1) -- (40,1);
\draw[thick,blue!60!black] plot file {figures/cascade_K4.dat};
\draw[thick,red!70!black] plot file {figures/cascade_K12.dat};
\node[font=\scriptsize,blue!60!black,anchor=west] at (16,0.55) {$K=4$: устанавливается};
\node[font=\scriptsize,red!70!black,anchor=west] at (16,1.75) {$K=12$: пульсирует};
\end{tikzpicture}
\caption{Каскад из трёх ступеней с обратной связью по~\eqref{eq:cascade}, уровень конечного гормона в долях равновесного $x_3^*$. При усилении петли ниже восьми уровень устанавливается; выше --- каскад сам порождает ровные пульсации.}
\label{fig:cascade}
\end{figure}

Настоящие гормоны каскадов действительно выходят пульсами: уровень гормона стресса, например, колеблется с периодом около часа. По одной из гипотез, эти пульсации порождает сама обратная связь между ступенями с запаздыванием, и отдельный генератор ритма для них не нужен~\cite{Walker2010}. Это та же причина колебаний, что и у петли \nt{ГЛУТ}--\nt{ГАМК} в утверждении~\ref{prop:ei} и у петли растяжения мышцы в условии~\eqref{eq:delaystable}: отрицательная обратная связь, которая запаздывает.

\section{Импульсный код гормонов}

Пульсы бывают не только побочным эффектом, но и кодом. Нейроны, управляющие размножением, выпускают свой сигнал в кровь короткими порциями примерно раз в час. Если железе-получателю подать тот же сигнал не порциями, а непрерывно, она не работает в полную силу, как при пульсах, а замолкает; если снова подавать порциями раз в час, работа восстанавливается~\cite{Belchetz1978}. Значит, получатель читает не уровень, а \emph{частоту} порций.

Для инженера здесь нет загадки. Рецептор, который при долгом действии вещества устаёт и перестаёт отвечать, --- фильтр верхних частот, как истощающийся синапс~\eqref{eq:depression} из главы~\ref{sec:code}. Непрерывный сигнал он глушит, а перемены пропускает. Такому получателю можно сообщить что-то только импульсами, и информация переходит из уровня в частоту и форму порций. Разные частоты порций к тому же по-разному действуют на разные клетки одной и той же железы, так что по одной химической линии можно передать больше одной величины --- как по одному проводу \nt{ДОФА} в главе~\ref{sec:dofaalt} шли быстрая и медленная составляющие.

\section{Суточный ритм: медленный генератор режимов}

Самый медленный ритм тела --- суточный. Его порождает отрицательная обратная связь с запаздыванием внутри клеток: гены производят белки, которые с задержкой в несколько часов подавляют собственное производство~\cite{TakahashiJS2017}. Опять запаздывающая отрицательная обратная связь --- четвёртый раз в этой книге, --- и опять ритм, на этот раз с периодом около суток. Главный генератор --- небольшая группа из десятков тысяч нейронов в основании мозга, ритм которых синхронизирует остальные клетки тела.

Собственный период этого генератора у человека не ровно сутки, а в среднем $24.18$ часа~\cite{Czeisler1999}. Каждый день его подводит свет, приходящий от датчиков глаза (глава~\ref{sec:sensors}). Для электронщика это \emph{фазовая автоподстройка частоты}: генератор со своей частотой $\omega_0$ подстраивается к внешнему сигналу с частотой $\omega$, и разность фаз $\varphi$ подчиняется уравнению
\begin{empheq}[box=\keybox]{equation}
\frac{\dd\varphi}{\dd t} \;=\; (\omega-\omega_0) \;-\; K_\varphi\,\sin\varphi .
\label{eq:pll}
\end{empheq}
Подстройка удерживается, если $|\omega-\omega_0|<K_\varphi$. Генератору с периодом $24.18$ часа нужно каждый день сдвигаться примерно на одиннадцать минут; утреннего света для такого сдвига обычно хватает. В полярную ночь или в пещере без света подстройки нет, и ритм уходит на собственный период.

Суточный генератор --- не тактовый генератор компьютера. Он не отсчитывает шаги вычислений и не синхронизирует импульсы нейронов. Он переключает \emph{режимы}: сон и бодрствование, уровни гормонов, температуру тела, готовность к еде. Это медленный широковещательный регистр того же рода, что и регистры глав~\ref{sec:serocomputer}--\ref{sec:acet}, только его значение меняется по расписанию, подведённому к солнцу.

\begin{keyblock}
\begin{proposition}[Химический выход]
\label{prop:chemoutput}
Медленный выход машины построен из тех же деталей, что и быстрый. Органы получают два провода противоположного знака и разной скорости --- газ \nt{НОРА} и тормоз \nt{АЦЕТ}, --- а весь организм сразу --- широковещательные сигналы по крови, в том числе адреналин \nt{АДРЕ}. Уровни гормонов держат каскады с отрицательной обратной связью, точность которых ограничена устойчивостью; часть гормонов кодируется частотой порций; суточный ритм --- медленный генератор с фазовой подстройкой по свету. Устройство путей и опыты с порциями и с периодом измерены; объяснение пульсаций каскада самой обратной связью --- гипотеза.
\end{proposition}
\end{keyblock}

\section{Итог}

\begin{table}[t]
\centering
\footnotesize
\setlength{\tabcolsep}{4pt}
\renewcommand{\arraystretch}{1.15}
\begin{tabular}{>{\raggedright}p{3.1cm}>{\raggedright}p{4.8cm}>{\raggedright\arraybackslash}p{4.9cm}}
\toprule
Что делает химический выход & Как & Что известно \\
\midrule
подстраивает органы & газ \nt{НОРА} и тормоз \nt{АЦЕТ}~\eqref{eq:gasbrake} & измерено; тормоз быстрее газа \\
переводит всё тело в режим бега & адреналин \nt{АДРЕ} в кровь & измерено \\
держит уровни гормонов & каскад с обратной связью, $K<8$~\eqref{eq:cascadestable} & каскады и обратная связь измерены; пульсации от самой петли --- гипотеза \\
передаёт частотой порций & уставший рецептор как фильтр верхних частот & зависимость от порций измерена \\
переключает режимы за сутки & генератор с фазовой подстройкой по свету~\eqref{eq:pll} & период и подстройка светом измерены \\
\bottomrule
\end{tabular}
\caption{Как машина управляет химией тела.}
\label{tab:chemoutput}
\end{table}

У машины два выхода (таблица~\ref{tab:chemoutput}). Быстрый --- мышцы: миллисекунды, адресно, в каждый сустав. Медленный --- химия тела: секунды, минуты и сутки, по двум проводам к органам и по крови ко всем сразу. Детали у обоих выходов одни и те же --- два провода противоположного знака, ритмоводители, обратная связь и её запаздывание, --- и те же, из которых собрана машина внутри.
